% GENERATED FILE. Edit canonical lessons, then run npm run generate:books.
% canonical-source-hash: fnv1a64:db4c8083c1fe50d2
% canonical-lessons: AR-C122-khizana, AR-C122-musht, AR-C122-sabun, AR-C122-minshafa, AR-C122-firash

\chapter{Cupboard, Comb, Soap, Towel, Mattress}
\label{ch:ar-cupboard-comb-soap-towel-mattress}
\hlchaptermodality{122}

\begin{chapteropening}
\textbf{By the end of this chapter:} \emph{I can say a cupboard, a comb, soap, a towel and a mattress.}

Complete the last lesson of chapter 122: I can say a cupboard, a comb, soap, a towel and a mattress.
\end{chapteropening}

\section[\texorpdfstring{\ar{خزانة} (khizāna)}{خزانة (khizāna)}]{\ar{خزانة} (khizāna) — a cupboard}

\label{lesson:AR-C122-khizana}



\begin{quote}
Before the new one: say the Arabic for thyme, then the Arabic for cumin.
\end{quote}

\subsection*{You'll want to know: \ar{خزانة}}
\textbf{\ar{خزانة}} — \emph{khizāna} — \textquotedblleft{}a cupboard\textquotedblright{}.

\begin{cousinweb}
\textbf{\ar{خزانة}} (\emph{khizāna}) is built on the native Arabic root \emph{kh-z-n}, \textquotedblleft{}to store, to lay up\textquotedblright{} (Lane). \emph{Khizāna} is a store-place; English \emph{magazine} comes from \emph{makhāzin}, \textquotedblleft{}storehouses\textquotedblright{}, on the same root.
\end{cousinweb}

\begin{grammarlens}[title={What you've built}]
One more word for this chapter.
\end{grammarlens}

\subsection*{Your turn}
\begin{itemize}
\raggedright
  \item \emph{Say it:} \emph{khizāna}
  \item \emph{Say it:} \emph{khizāna}, once more
  \item \emph{Say it:} say \emph{khizāna}
  \item \emph{Recall:} say the Arabic for thyme, then the Arabic for cumin, then say \emph{khizāna} again
\end{itemize}

\begin{tcolorbox}[breakable,colback=teal!4,colframe=teal!35!black,title={Before you move on}]
What does \ar{خزانة} mean? (\textquotedblleft{}A cupboard\textquotedblright{}.) Say it once more.
\end{tcolorbox}
\section[\texorpdfstring{\ar{مشط} (mushṭ)}{مشط (mushṭ)}]{\ar{مشط} (mushṭ) — a comb}

\label{lesson:AR-C122-musht}



\begin{quote}
Before the new one: say the Arabic for cumin, then the Arabic for a cupboard.
\end{quote}

\subsection*{You'll want to know: \ar{مشط}}
\textbf{\ar{مشط}} — \emph{mushṭ} — \textquotedblleft{}a comb\textquotedblright{}.

\begin{cousinweb}
\textbf{\ar{مشط}} (\emph{mushṭ}) is built on the native Arabic root \emph{m-sh-ṭ}, \textquotedblleft{}to comb\textquotedblright{} (Lane). \emph{Mushṭ} is a comb.
\end{cousinweb}

\begin{grammarlens}[title={What you've built}]
One more word for this chapter.
\end{grammarlens}

\subsection*{Your turn}
\begin{itemize}
\raggedright
  \item \emph{Say it:} \emph{mushṭ}
  \item \emph{Say it:} \emph{mushṭ}, once more
  \item \emph{Say it:} say \emph{mushṭ}
  \item \emph{Recall:} say the Arabic for cumin, then the Arabic for a cupboard, then say \emph{mushṭ} again
\end{itemize}

\begin{tcolorbox}[breakable,colback=teal!4,colframe=teal!35!black,title={Before you move on}]
What does \ar{مشط} mean? (\textquotedblleft{}A comb\textquotedblright{}.) Say it once more.
\end{tcolorbox}
\section[\texorpdfstring{\ar{صابون} (ṣābūn)}{صابون (ṣābūn)}]{\ar{صابون} (ṣābūn) — soap}

\label{lesson:AR-C122-sabun}



\begin{quote}
Before the new one: say the Arabic for a cupboard, then the Arabic for a comb.

Say \textquotedblleft{}green\textquotedblright{} and \textquotedblleft{}yellow.\textquotedblright{} (\emph{Akhḍar}, the root of \textquotedblleft{}vegetables\textquotedblright{}; \emph{aṣfar}, the root of \emph{ṣifr}, \textquotedblleft{}zero.\textquotedblright{})
\end{quote}

\subsection*{You'll want to know: \ar{صابون}}
\textbf{\ar{صابون}} — \emph{ṣābūn} — \textquotedblleft{}soap\textquotedblright{}.

\begin{cousinweb}
\textbf{\ar{صابون}} (\emph{ṣābūn}) is usually traced to Latin \emph{sāpō}, \textquotedblleft{}soap\textquotedblright{} (EtymDB), probably through Aramaic or Greek. Lane quotes a grammarian who derived it from an Arabic verb, a folk etymology, and another who called it a word common to many languages.
\end{cousinweb}

\begin{grammarlens}[title={What you've built}]
One more word for this chapter.
\end{grammarlens}

\subsection*{Your turn}
\begin{itemize}
\raggedright
  \item \emph{Say it:} \emph{ṣābūn}
  \item \emph{Say it:} \emph{ṣābūn}, once more
  \item \emph{Say it:} say \emph{ṣābūn}
  \item \emph{Recall:} say the Arabic for a cupboard, then the Arabic for a comb, then say \emph{ṣābūn} again
\end{itemize}

\begin{tcolorbox}[breakable,colback=teal!4,colframe=teal!35!black,title={Before you move on}]
What does \ar{صابون} mean? (\textquotedblleft{}Soap\textquotedblright{}.) Say it once more.
\end{tcolorbox}
\section[\texorpdfstring{\ar{منشفة} (minshafa)}{منشفة (minshafa)}]{\ar{منشفة} (minshafa) — a towel}

\label{lesson:AR-C122-minshafa}



\begin{quote}
Before the new one: say the Arabic for a comb, then the Arabic for soap.
\end{quote}

\subsection*{You'll want to know: \ar{منشفة}}
\textbf{\ar{منشفة}} — \emph{minshafa} — \textquotedblleft{}a towel\textquotedblright{}.

\begin{cousinweb}
\textbf{\ar{منشفة}} (\emph{minshafa}) is built on the native Arabic root \emph{n-sh-f}, \textquotedblleft{}to dry, to wipe dry\textquotedblright{} (Lane). \emph{Minshafa} is an instrument noun, \textquotedblleft{}a thing for drying\textquotedblright{}.
\end{cousinweb}

\begin{grammarlens}[title={What you've built}]
One more word for this chapter.
\end{grammarlens}

\subsection*{Your turn}
\begin{itemize}
\raggedright
  \item \emph{Say it:} \emph{minshafa}
  \item \emph{Say it:} \emph{minshafa}, once more
  \item \emph{Say it:} say \emph{minshafa}
  \item \emph{Recall:} say the Arabic for a comb, then the Arabic for soap, then say \emph{minshafa} again
\end{itemize}

\begin{tcolorbox}[breakable,colback=teal!4,colframe=teal!35!black,title={Before you move on}]
What does \ar{منشفة} mean? (\textquotedblleft{}A towel\textquotedblright{}.) Say it once more.
\end{tcolorbox}
\section[\texorpdfstring{\ar{فراش} (firāsh)}{فراش (firāsh)}]{\ar{فراش} (firāsh) — a mattress}

\label{lesson:AR-C122-firash}



\begin{quote}
Before the new one: say the Arabic for soap, then the Arabic for a towel.
\end{quote}

\subsection*{You'll want to know: \ar{فراش}}
\textbf{\ar{فراش}} — \emph{firāsh} — \textquotedblleft{}a mattress\textquotedblright{}.

\begin{cousinweb}
\textbf{\ar{فراش}} (\emph{firāsh}) is built on the native Arabic root \emph{f-r-sh}, \textquotedblleft{}to spread out\textquotedblright{} (Lane). \emph{Firāsh} is a thing spread to lie on.
\end{cousinweb}

\begin{grammarlens}[title={What you've built}]
One more word for this chapter.
\end{grammarlens}

\subsection*{Your turn}
\begin{itemize}
\raggedright
  \item \emph{Say it:} \emph{firāsh}
  \item \emph{Say it:} \emph{firāsh}, once more
  \item \emph{Say it:} say \emph{firāsh}
  \item \emph{Recall:} say the Arabic for soap, then the Arabic for a towel, then say \emph{firāsh} again
\end{itemize}

\begin{tcolorbox}[breakable,colback=teal!4,colframe=teal!35!black,title={Before you move on}]
What does \ar{فراش} mean? (\textquotedblleft{}A mattress\textquotedblright{}.) Say it once more.
\end{tcolorbox}
